\chapter{Considerações Finais}
\label{cap:consideracoes}

Este relatório apresentou, em conformidade com os requisitos da Etapa 1 da disciplina de Banco de Dados, a análise detalhada do problema de monitoramento e manutenção colaborativa de infraestrutura urbana nos municípios brasileiros, o seu enquadramento fundamentado nos Objetivos de Desenvolvimento Sustentável (ODS) da Agenda 2030 da Organização das Nações Unidas e a delimitação do escopo conceitual do sistema de banco de dados proposto.

A investigação do contexto nacional, respaldada pelas estatísticas oficiais do Censo Demográfico 2022 \cite{ibge2024censo2022malha, ibge2025entorno} e pelos diagnósticos do Sistema Nacional de Informações sobre Saneamento \cite{snis2023drenagem, snis2023residuos}, evidenciou que a conservação das vias e equipamentos coletivos afeta mais de 174 milhões de moradores urbanos. Constatou-se que os entraves enfrentados pelas prefeituras decorrem sobretudo da fragmentação informacional entre os canais de recepção de chamados e a retaguarda operacional das secretarias municipais, gerando protocolos duplicados, imprecisão geográfica, morosidade no despacho de equipes e frotas, descontrole sobre o consumo de insumos de almoxarifado e ausência de retorno formal ao cidadão solicitante.

Sob a perspectiva da sustentabilidade global e das políticas públicas urbanas, demonstrou-se que a estruturação de um banco de dados relacional para gerir esse ciclo completo responde diretamente ao \textbf{ODS 11 (Cidades e Comunidades Sustentáveis)} \cite{onubrasil2015ods11, ipea2019ods11}, com ênfase nas \textbf{Metas 11.3, 11.6 e 11.7}. De modo complementar, o projeto fortalece a resiliência da infraestrutura municipal preconizada pelo \textbf{ODS 9 (Meta 9.1)} \cite{onubrasil2015ods9} e a transparência e eficácia institucional exigidas pelo \textbf{ODS 16 (Metas 16.6 e 16.7)} \cite{onubrasil2015ods16}, alinhando-se às diretrizes do Estatuto da Cidade \cite{brasil2001estatutocidade} e do Código de Defesa dos Usuários dos Serviços Públicos \cite{brasil2017lei13460}.

Por fim, a definição do escopo do banco de dados organizou o minimundo em três subdomínios coesos --- demandas do cidadão e base geográfica, estrutura operacional municipal e execução de ordens de serviço articulada à cadeia de suprimentos ---, totalizando doze entidades de negócio e duas relações associativas de controle de estoque e compras por lote. Como passo seguinte, o trabalho articula-se diretamente à \textbf{Etapa 2}, dedicada à formalização do Modelo Entidade-Relacionamento (MER) conceitual na notação do brModelo (doze entidades, setenta atributos de entidades, quatorze relacionamentos binários e seis atributos de relacionamentos), seguida do mapeamento de \citeonline{elmasri2019sistemas} para as quatorze tabelas do esquema relacional e sua validação em SQL DDL.
